\section{Kết luận}
\label{sec:conclusion}

\subsection{Trả lời câu hỏi nghiên cứu}

\begin{description}[leftmargin=2.6em,labelindent=0pt]
  \item[RQ1.] Ở mức nén $\approx$4$\times$, \textbf{chọn lọc theo câu} (Provence $\tau=0{,}1$, 4,6$\times$)
  là cách duy nhất không làm giảm accuracy có ý nghĩa thống kê ở cả ba vị trí (64,0\% so với 65,3\%
  của ngữ cảnh đầy đủ). Chọn lọc theo tài liệu (rerank top-5) rẻ hơn nhiều nhưng mất 2,6 điểm và giảm
  có ý nghĩa khi đoạn vàng ở cuối. Xoá theo token (LLMLingua-2 r=0,25) làm accuracy rơi về mức
  closed-book (38,3\% so với 38,8\%): phân tích lỗi cho thấy phần lớn thiệt hại đến từ văn bản bị xé
  vụn chứ không phải mất đáp án. Với reader hiện đại, đường cong vị trí vốn đã gần phẳng (chênh 5,8
  điểm) nên nén không còn tác dụng ``làm phẳng'' như báo cáo cho GPT-3.5.
  \item[RQ2.] Ở quy mô 50,7k đoạn, retriever dày đặc nhỏ (bge-small, 33M tham số) vượt BM25 21 điểm
  hit@5; lượng tử int8 giảm chỉ mục 4$\times$ không mất chất lượng, nhị phân + rescore giảm 32$\times$
  chỉ mất 2,8 điểm.
  \item[RQ3.] Tái lập được hai kết luận của các công bố: BM25 kém xa dense, và LLMLingua-2 dùng đơn lẻ
  ở 4$\times$ sụt mạnh ($-$21 đến $-$30 điểm so với $-$11 đến $-$27 được công bố). Không tái lập được
  kết luận ``nén làm tăng accuracy khi đoạn vàng ở giữa'' của LongLLMLingua ($+17{,}1$ điểm với GPT-3.5,
  $+1{,}0$ không có ý nghĩa với Qwen2.5-7B): hiệu ứng đó phụ thuộc vào reader bị \emph{lost in the
  middle}, điều không còn đúng với LLM mở thế hệ mới ở độ dài 3K token.
  \item[RQ4.] Có. Bộ nén chạy trên CPU (rerank 0,2 s/câu, Provence 4,9 s/câu), reader gọi qua API;
  nén giảm 78\% token đầu vào và do đó giảm chi phí API tương ứng. GPU chỉ cần để rút ngắn độ trễ
  nén.
  \item[RQ5.] % TODO(T18): fill with PISCO / own soft model numbers
\end{description}

\subsection{Hạn chế}

\begin{itemize}
  \item Chỉ một bộ dữ liệu (NQ-Open, câu trả lời ngắn) và một reader chính; kết luận về nén mức token
  có thể khác với tác vụ cần lập luận nhiều bước hoặc ngữ cảnh dài hơn nhiều (32K+).
  \item Chỉ số \emph{best subspan EM} phạt các câu trả lời đúng nhưng khác định dạng (ví dụ ngày
  tháng); một phần lỗi trong Mục~\ref{sec:res-errors} thuộc loại này.
  \item LongLLMLingua chỉ chạy 200 câu ở một vị trí vì chi phí; các số liệu GPT-3.5 chỉ dùng để so
  sánh độ chênh, không so sánh tuyệt đối.
  \item Mô hình nén mềm của nhóm được huấn luyện 1 epoch trên 15k mẫu với một GPU 12GB, nhỏ hơn nhiều
  so với PISCO (huấn luyện trên hàng trăm nghìn mẫu).
\end{itemize}

\subsection{Hướng phát triển}

\begin{itemize}
  \item Kết hợp hai nhánh hiệu quả nhất: chọn lọc theo câu (Provence) cho tài liệu được xếp hạng cao
  và bộ nhớ nén cho phần còn lại -- tương tự cơ chế trộn token chọn lọc và token nén trong attention
  thưa của DeepSeek.
  \item Đánh giá trên ngữ cảnh dài hơn (50--100 tài liệu) và các bộ dữ liệu đa bước (HotpotQA,
  MuSiQue), nơi reader còn bị ảnh hưởng bởi vị trí.
  \item Chưng cất Provence sang mô hình nhỏ hơn để giảm độ trễ trên CPU xuống dưới 1 s/câu.
\end{itemize}
